\section{Evidence and measure: the charge with its witnesses, and the trace that measures it}

A bare reading (\cref{thm:charge}) is the presence pattern alone; positing
does not admit one directly. What is admitted carries its witnesses:
\texttt{Valued} pairs a value with the lossless evidence it was read off,
and the trace-level \texttt{measure} builds that evidence from an
articulating trace against a codec's contradiction, never by a denial primitive.

\begin{proposition}[Charged: value with lossless evidence]\label{thm:charged}\lean{Obsmetrologia.Valued}\leanok\uses{thm:charge}
A \emph{charged} value pairs a \texttt{value : Charge} with
\texttt{evidence : Slot \(\times\) Slot} (for-witnesses, against-witnesses).
\texttt{presence} reads the evidence pair back as a reading. Evidence belongs
to value, which enters at positing: an articulation cell's subject and object
slots are positions of a predication, not evidence, and are not read as a
value. \texttt{joinK}, \texttt{meetK}, \texttt{joinT},
\texttt{meetT} apply the matching \texttt{Charge} operation to the value
while unioning the evidence on each side regardless of which operation is
applied --- combining evidence is never selective. \texttt{negate} swaps the
two evidence sides along with the value; \texttt{conflate} may raise the
value above what the evidence shows, the one operation licensed to do so.
\end{proposition}

\begin{proposition}[Measured]\label{thm:charged-measured}\lean{Obsmetrologia.Valued.Measured}\leanok\uses{thm:charged}
A charged value is \emph{measured} when its value is exactly its evidence's
presence pattern: \(x.\mathrm{value} = x.\mathrm{presence}\). This is the
formal invariant behind ``measured, not seeded'': a measured value carries no
more and no less than its witnesses show.
\end{proposition}

\begin{theorem}[A one-sided value is measured]\label{thm:valued-onesided-measured}\lean{Obsmetrologia.Valued.oneSided_measured}\leanok\uses{thm:charged-measured}
A value witnessed on its first coordinate only, by one token, is measured: its
value is exactly the presence pattern its single witness shows.
\end{theorem}

\begin{theorem}[Knowledge-join preserves measured]\label{thm:charged-joink-measured}\lean{Obsmetrologia.Valued.joinK_measured}\leanok\uses{thm:charged-measured}
If \(x\) and \(y\) are both measured, so is \(\texttt{joinK}\,x\,y\): raising
two honestly-read values by their honest union of evidence keeps the result
honest.
\end{theorem}

\begin{theorem}[Absurd reached, witnesses kept: the meet can break measured]\label{thm:charged-meetk-not-measured}\lean{Obsmetrologia.Valued.meetK_not_measured}\leanok\uses{thm:charged-measured, thm:charged-joink-measured}
There exist measured \(x\) and \(y\) with \(\texttt{meetK}\,x\,y\)
\textbf{not} measured: the knowledge meet of a potential telling and an
active telling drives the value down to \(\bot_k\) while keeping both
witnesses on record. This is the formal difference between ``absurd, no
evidence at all'' and
``absurd, two tellings that cancelled'' (evidenced on both sides, value below
what the evidence alone would license reading as neutral) --- the meet
reaches a value the evidence has outrun, not a value the evidence supports.
\end{theorem}

\begin{theorem}[Presumption only ever raises]\label{thm:charged-presumefree-ge}\lean{Obsmetrologia.Valued.presumeFree_ge}\leanok\uses{thm:charged}
\texttt{presumeFree} raises every evidence-free side of a charged value to
present while leaving the evidence itself untouched. For any \(x\),
\(x.\mathrm{value} \le_k \texttt{presumeFree}\,x.\mathrm{value}\): presumption
is the one sanctioned breach of \cref{thm:charged-measured}, and it never
lowers a value, only raises one that evidence has not yet pinned.
\end{theorem}

\begin{theorem}[A one-sided value presumes to stable]\label{thm:charged-presumefree-articulate}\lean{Obsmetrologia.Valued.presumeFree_oneSided}\leanok\uses{thm:charged-presumefree-ge}
Presuming a value witnessed on one coordinate only reaches \(\top_k\):
potential presumes up to stable. This is the pointwise fact the
counit cone (\cref{thm:cvaluation-brim}) generalizes to every cell of a
valuation.
\end{theorem}

\begin{definition}[Contradiction]\label{thm:contrary}\lean{Obsmetrologia.Contradictory}\leanok\uses{def:claim}
A \emph{contradiction} is an involution \(c : \texttt{Claim} \to
\texttt{Claim}\) on claims, together with the proof \(c\,(c\,p) = p\) for
every \(p\). It pairs each claim with the mark of its denial, told as a mark
in its own right, so that evidence against \(\varphi\) is evidence for its
contradictory at the same subject and object: denial is a telling against the
same atom, as in Belnap's four-valued logic (Belnap 1977; Dunn 1976), and no
denial operation is assumed. Being an involution, it is contradiction rather
than contrariety: contraries such as \emph{likes} and \emph{dislikes} can both
be false (Horn 1989), so a contrary is a separate predicate related to its
partner by a rule settling applies, not by this pairing.
\end{definition}

\begin{definition}[Witnesses of a claim in a trace]\label{thm:witnesses}\lean{Obsmetrologia.witnesses}\leanok\uses{thm:oriented-cell, thm:predication}
\(\texttt{witnesses}\,\psi\,\texttt{trace}\) collects, by explicit structural
recursion over a list of articulatings, the head token of every articulating
in the trace whose oriented cell (\cref{thm:oriented-cell}) equals \(\psi\).
It is a \texttt{Slot}: a list, built without \texttt{filter} or \texttt{map},
so that every lemma about it is a structural induction.
\end{definition}

\begin{definition}[Measuring a trace]\label{thm:trace-measure}\lean{Obsmetrologia.measure}\leanok\uses{thm:witnesses, thm:contrary, thm:charged}
\(\texttt{measure}\,k\,\texttt{trace}\,\varphi\) reads a \texttt{Charged}
value at the predication cell \(\varphi\) from a trace of articulatings and a
contradiction \(k\): for-witnesses are \(\texttt{witnesses}\,\varphi\,
\texttt{trace}\), against-witnesses are
\(\texttt{witnesses}\,(\texttt{contradictory}\,k\,\varphi)\,\texttt{trace}\), and
the value is the presence pattern of that evidence pair. There is no
operation asserting a charge directly; a claim is denied only by its
contradictory having been articulated somewhere in the trace.
\end{definition}

\begin{theorem}[Measuring a trace is always measured]\label{thm:measure-measured}\lean{Obsmetrologia.measure_measured}\leanok\uses{thm:trace-measure, thm:charged-measured}
\(\texttt{measure}\,k\,\texttt{trace}\,\varphi\) is measured for every \(k\),
\texttt{trace}, and \(\varphi\): the value is built as the presence pattern
of exactly the evidence \texttt{witnesses} collects, by construction.
\end{theorem}

\begin{theorem}[The empty trace tells nothing]\label{thm:measure-nil}\lean{Obsmetrologia.measure_nil}\leanok\uses{thm:trace-measure, thm:four-charges}
Measuring the empty trace gives \(\bot_k\) at every claim: no articulating
means no witnesses on either side, so the reading is the knowledge-least
value by construction, not by a seeded default.
\end{theorem}

\begin{theorem}[Evidence never retracts]\label{thm:measure-cons-le}\lean{Obsmetrologia.measure_cons_le}\leanok\uses{thm:trace-measure, thm:kle}
Extending a trace by one more articulating can only raise the measured value
in \(\le_k\), never lower it: \(\texttt{measure}\,k\,\texttt{trace}\,\varphi
\le_k \texttt{measure}\,k\,(a :: \texttt{trace})\,\varphi\). This is time's
arrow at the level of measured evidence: knowledge is joined, never
withdrawn.
\end{theorem}

\begin{theorem}[An articulating tells its own cell true]\label{thm:measure-self-pos}\lean{Obsmetrologia.measure_self_pos}\leanok\uses{thm:trace-measure, thm:oriented-cell}
Measuring the singleton trace \([a]\) at \(a\)'s own oriented cell gives a
positive reading: an articulating is always a for-witness of the cell it
itself generates.
\end{theorem}

\begin{theorem}[No denial]\label{thm:no-denial}\lean{Obsmetrologia.no_denial}\leanok\uses{thm:trace-measure, thm:contrary}
If measuring a trace at \(\varphi\) reads negative evidence, the
contradictory claim was itself articulated somewhere in the trace: negative
evidence exists only where the codec's contradictory claim was actually told,
never as a bare denial of \(\varphi\).
\end{theorem}

\begin{theorem}[Negation is measured]\label{thm:measure-contradictory-negate}\lean{Obsmetrologia.measure_contradictory_negate}\leanok\uses{thm:trace-measure, thm:contrary}
Measuring a trace at the contradictory of \(\varphi\) is the Belnap negation
of measuring it at \(\varphi\): the two evidence sides exchange. Negation on
readings is therefore derived from contradiction on marks, not assumed.
\end{theorem}

\begin{definition}[The measured composite of the spine]\label{thm:spine-runtrace}\lean{Obsmetrologia.Spine.runTrace}\leanok\uses{thm:trace-measure, thm:settling}
\(\texttt{runTrace}\,\texttt{closure}\,k\,\texttt{trace}\,\varphi\) is
articulate, predicate, \emph{measure}, settle: it settles the closure applied
to the value of \(\texttt{measure}\,k\,\texttt{trace}\,\varphi\) at the locus
\(\varphi\). Where \texttt{Spine.runSupplied} admits a caller-supplied
reading and constructs no admission of evidence, \texttt{runTrace} is the
composite that actually measures.
\end{definition}
\begin{definition}[The spine]\label{thm:spine-run}\lean{Obsmetrologia.Spine.run, Obsmetrologia.Spine.run_eq_runTrace, Obsmetrologia.Spine.run_reading_fixed}\leanok\uses{thm:spine-runtrace, thm:spine-predicate-oriented}
\(\texttt{run}\,\texttt{closure}\,k\,\texttt{trace}\,a\) is the spine:
articulate \(a\) into the trace, predicate its oriented cell, measure the
reading there from \(a :: \texttt{trace}\), and settle; that is,
\(\texttt{runTrace}\,\texttt{closure}\,k\,(a :: \texttt{trace})\,(\texttt{predicate}\,a)\).
No reading enters from the caller.
\end{definition}
\begin{theorem}[What is articulated is told for]\label{thm:spine-run-pos}\lean{Obsmetrologia.Spine.run_pos, Obsmetrologia.measure_cons_self_pos, Obsmetrologia.Spine.KLE_pos}\leanok\uses{thm:spine-run, thm:measure-self-pos}
The spine's settled reading at the articulated claim has its positive side
present, for every closure, contradiction and trace: the articulating is its
own for-witness in any trace it heads, and settling, being knowledge-inflationary,
never withdraws told-for evidence.
\end{theorem}

\begin{theorem}[The empty trace's spine is the knowledge-least valuation]\label{thm:spine-measure-nil-knowledgeleast}\lean{Obsmetrologia.Spine.measure_nil_knowledgeLeast}\leanok\uses{thm:measure-nil, thm:knowledge-least}
For every claim \(\varphi\), measuring the empty trace and reading it as a
valuation is exactly \texttt{Valuation.knowledgeLeast}~\(\varphi\). This
anchors \cref{thm:knowledge-least}'s own reading --- the value at cells no
articulating has reached --- as the special case of measuring nothing at all.
\end{theorem}

\begin{definition}[The measured Chu space of a trace]\label{thm:measured-chu}\lean{Obsmetrologia.measuredChu}\leanok\uses{thm:trace-measure}
A trace's measurement read as a FOUR-valued Chu space: tokens are the points, a
claim together with the token it is told toward is a test, and the entry at a
token \(t\) and test \((c, o)\) is the reading measured at the oriented cell
\(\langle t, c, o\rangle\).
\end{definition}
\begin{definition}[The root identification: biextensional collapse]\label{thm:root-identification}\lean{Obsmetrologia.ChuSpace.collapse, Obsmetrologia.ChuSpace.collapse_separable, Obsmetrologia.ChuSpace.collapse_extensional, Obsmetrologia.ChuSpace.collapse_mk_eq_iff}\leanok\uses{thm:measured-chu}
The biextensional collapse of a Chu space identifies points with equal rows and
states with equal columns, and the matrix descends to both quotients. The result
is separable and extensional (Pratt 1999; clarification of a formal context,
Ganter \& Wille 1999; the separated, extensional classification, Barwise \&
Seligman 1997). Identifying two presentations uses \texttt{Quot.sound} and
nothing else.
\end{definition}
\begin{theorem}[Measurement factors through the root identification]\label{thm:measure-factors-collapse}\lean{Obsmetrologia.measure_factors_collapse, Obsmetrologia.token_identified_iff}\leanok\uses{thm:root-identification}
The reading measured at any cell is the collapsed matrix at the cell's collapsed
subject and collapsed test, with no axiom. Two tokens are identified exactly when
every claim toward every token measures them alike: identity is role.
\end{theorem}
\begin{theorem}[Chu transforms descend]\label{thm:chu-collapse-descends}\lean{Obsmetrologia.ChuMap.collapse}\leanok\uses{thm:root-identification}
A Chu transform \(C \to D\) induces a Chu transform between the collapses: its
forward map respects row equality and its backward map column equality, both by
adjointness.
\end{theorem}
